\section{Case E Continuous dependent sensing}\label{sec:caseE}

\subsection{Total dependent Borel kernels}
Work in the standard-Borel part of the QBS probability interpretation
of O3~\cite{O3}. Define the Borel subspaces
\[
A=(0,1],\quad E=\{(t,y):t\in A,\ 0\leq y\leq t\},\quad
F=\{(t,y,z):(t,y)\in E,\ 0\leq z\leq y\}.
\]
The projection interfaces $E\to A$ and $F\to E$ are deterministic;
sampling is a computation, not a deterministic interface map. Let
$U,V$ be independent uniform random variables on $[0,1]$ and define
\[
K(t)=\Law(t,tU),\qquad L(t,y)=\Law(t,y,yV).
\]
The continuation is typed at the actual observed $(t,y)$ and uses the
result-dependent interval $[0,y]$. At $y=0$ it is the Dirac law at
$(t,0,0)$, so it is defined on the entire dependent domain, not merely
almost surely.

The realizing maps $(t,u)\mapsto(t,tu)$ and
$(t,y,v)\mapsto(t,y,yv)$ are continuous on their domains. Integration
of a bounded Borel function against the fixed uniform source is Borel
in its parameter, proving kernel measurability. Their Kleisli composite
is explicitly
\begin{equation}\label{eq:E-bind}
(L\star K)(t)=\Law(t,tU,tUV),\qquad
\mathbb E[h]=\int_0^1\int_0^1 h(t,tu,tuv)\,dv\,du
\end{equation}
for every Borel $h:F\to[0,1]$. This is forward-generative bind; it uses
no chosen disintegration or general measurable-coupling theorem. An
arbitrary input law $\mu$ gives the pushforward of
$\mu\otimes\operatorname{Unif}\otimes\operatorname{Unif}$ under the
same map, preserving the dependence between $t,y,z$.

\subsection{A nontrivial observed result}
The Borel test $O(t,y,z)=\mathbf1_{z\leq t/2}$ satisfies, for every
$t>0$,
\begin{equation}\label{eq:E-observation}
\Pr[O=1]=\int_0^{1/2}1\,du+
\int_{1/2}^1\frac{1}{2u}\,du=\frac{1+\log2}{2}.
\end{equation}
For $0<u\leq1/2$ all $v$ pass; for $u>1/2$ precisely the interval
$[0,1/(2u)]$ passes. The $u=0$ row is typed correctly although null.
More generally $\Pr[z\leq at]=a(1-\log a)$ for $0<a\leq1$, with
value zero at $a=0$. Thus $z/t=UV$ is not uniform, whereas
$\Pr[y\leq t/2]=1/2$. The two-stage dependent continuation has a
calculable effect, not just a continuous parameter name.

\begin{proposition}[Exact finite-history boundary]
No normalized finite family of deterministic raw histories followed by
a deterministic output map realizes $K(t)$ for any $t>0$. The same
obstruction holds for the composed $z$ marginal.
\end{proposition}
\begin{proof}
A finite raw family has finite image $S$ and assigns it mass one.
The uniform law on $[0,t]$ assigns every finite set mass zero.
These laws cannot agree. For the composite, conditional on any $u>0$,
$tuV$ has no atoms; the exceptional $u=0$ has measure zero.
Integrating shows that its $z$ marginal is also atomless, giving the
same finite-image contradiction.
\end{proof}
The obstruction concerns exact finite-support realization, not
approximation, infinite-history models or a finite program allowed an
atomless primitive. This case supplies an explicit two-step dependent
QBS/Borel protocol and observation. It does not claim an O1 finite
augmented-family realization, O4 finite-language interpretation, O5
finite local-coupling theorem, general continuous higher-order adequacy,
or a measurable selection of couplings. No output-dependent number of
holes is used: dependence is in the actual continuation fibre.

